\originalchapter{正则系综与配分函数}\label{ch:canonical}
\originalsection{从热库到Boltzmann权重}
把研究系统与远大于它的热库组成孤立整体, 两者相互作用能可忽略. 系统处于微观态$i$, 能量$E_i$时, 热库可用状态数为$\Omega_b(E_{\rm tot}-E_i)$. 展开热库熵得
\[
S_b(E_{\rm tot}-E_i)=S_b(E_{\rm tot})-\frac{E_i}{T}
-\frac{E_i^2}{2T^2 C_b}+\cdots.
\]
热库热容$C_b$足够大时, 一阶项支配微观能量差, 因而$p_i=\e^{-\beta E_i}/Z$, $\beta=1/(\kb T)$, $Z=\sum_i\e^{-\beta E_i}$. 有限热库修正显示适用条件, 不能把任意小热库都当作恒温器. 经典系统把求和换为$Z=\int\e^{-\beta H}\dd\Gamma$, 采用与状态计数一致的无量纲测度.

权重属于\textbf{微观状态}. 若问能量密度, 还要乘态密度: $p_E(E)=D(E)\e^{-\beta E}/Z$. 忽略这个因子就无法解释高温气体最概然能量非零. 简并能级$E$有$g(E)$个态, 总概率是$g(E)\e^{-\beta E}/Z$, 每个态的概率仍是$\e^{-\beta E}/Z$.

对$S=-\kb\sum_i p_i\ln p_i$代入权重, 得$S=\kb\ln Z+U/T$, 所以$F=U-TS=-\kb T\ln Z$. 若Hamilton量和测度不显含$T$, 固定所有外参求导可得
\[
U=-\partial_\beta\ln Z,\quad
\Var E=\partial_\beta^2\ln Z=\kb T^2 C_V,\quad
P=\kb T\partial_V\ln Z.
\]
对于有效的温度依赖Hamilton量, $-\partial_\beta\ln Z$可能还含$\beta\partial_\beta H$, 必须另核物理内能定义. 对通常温度无关模型, 热容非负是能量方差非负的直接结果.

\originalsection{独立子系统与经典理想气体}
若Hamilton量为$H_A+H_B$且相空间约束也分别独立, $Z=Z_AZ_B$, 自由能、能量与熵可加. $N$个固定位置上的相同振子给$Z=z^N$; $N$个同种自由经典粒子给$Z=z^N/N!$, 两者不可混淆.

单粒子动量Gaussian积分为$(2\pi m\kb T)^{3/2}$, 位置积分$V$, 因而$z=V/\lambda_T^3$, $\lambda_T=h/\sqrt{2\pi m\kb T}$是热de Broglie波长. 大$N$下
\[
\begin{aligned}
Z_N&=\frac1{N!}\left(\frac V{\lambda_T^3}\right)^N,\\
F&=-N\kb T\left[\ln\frac V{N\lambda_T^3}+1\right],\\
S&=N\kb\left[\ln\frac V{N\lambda_T^3}+\frac52\right],\quad
U=\frac32N\kb T.
\end{aligned}
\]
这与微正则结果接上. 经典稀薄条件为$n\lambda_T^3\ll1$, $n=N/V$. 当平均粒子间距与波长相当时, 只除$N!$仍不足以代替量子交换统计. 绝热路径由$\Delta S=0$给$VT^{3/2}=\text{常数}$.

固定振子$H=p^2/(2m)+Kq^2/2$有$z=\kb T/(\hbar\omega)$, $\omega=\sqrt{K/m}$. $p$和$q$在正则系综中独立, 各自方差为$m\kb T$和$\kb T/K$, 平均能量$\kb T$. 经典振子低温仍给热容$\kb$, 与第三律冲突, 提醒我们经典近似在量子能隙附近失效. 若频率随外参改变, 等熵条件是$T/\omega$固定; 只有固定频率时等熵才等温.

\originalsection{能量均分与两能级例子}
对一个取值覆盖整条实线的正则变量$x$, 在边界项$x\e^{-\beta H}$消失且积分收敛时, 分部积分给
\[
0=\int\partial_x(x\e^{-\beta H})\dd\Gamma
\quad\Longrightarrow\quad\avg{x\partial_xH}=\kb T.
\]
二次项$a x^2/2$因此贡献$\kb T/2$. 有硬边界的坐标或量子自由度要重新查条件, 不能把均分机械施加到所有自由度.

基态一个, 激发能$\epsilon$简并$g$重的单元, 有$z=1+g\e^{-\beta\epsilon}$, $p_{\rm exc}=g\e^{-\beta\epsilon}/z$. 若$N$个单元固定在不同位置, $U=N\epsilon p_{\rm exc}$,
\[
C=N\kb(\beta\epsilon)^2\frac{g\e^{-\beta\epsilon}}{(1+g\e^{-\beta\epsilon})^2},\quad
S=N\kb\left[\ln z+\beta\epsilon p_{\rm exc}\right].
\]
高温时激发概率趋$g/(g+1)$, 不是所有粒子都激发. 若能隙满足$\epsilon(V)=\epsilon_0(V/V_0)^{-\gamma}$, 固定温度求$P=-F_V$得$P=\gamma U/V$. 外参改变能谱, 因而即使固定晶格位置的内部激发也可以贡献压强.
